\section{From a one-sided prime-sum bound to zero-freeness}
\label{sec:window-analysis}

Let \(F=(\chi_i,a_i)_{i\in I}\) be a finite family of Dirichlet
characters, possibly of different positive moduli. Each weight \(a_i\)
is strictly positive. Each principal character in the family must have
modulus one. Write \(\Lambda(n)\) for the von Mangoldt function and
\(\gamma\) for Euler's constant. For \(y\geq1\) put
\[
 B_\chi(y)=\sum_{1\leq n\leq y}\Lambda(n)\chi(n)(y/n-1).
\]
Define
\[
 G_\chi(y)=
 \begin{cases}
 B_\chi(y)-y\log y+(1+\gamma)y,&\chi\text{ principal},\\
 B_\chi(y)+yL'(1,\chi)/L(1,\chi),&\chi\text{ nonprincipal}.
 \end{cases}
 \qquad G_F(y)=\sum_{i\in I}a_i\Re G_{\chi_i}(y).
\]
The second expression uses classical nonvanishing of \(L(1,\chi)\).
For definiteness, a zero of a meromorphic function is a point where its
order is positive. Consequently the pole of \(\zeta\) at one is excluded.
The formal library assigns a nonzero value to its total function at that
point, which gives the same assertion about zeros.

\paragraph{Window criterion.}
Fix \(0<\theta<1\), \(\delta>0\), and \(C\in\mathbb R\). Suppose that
for every \(Y_0\) there is \(Y\geq\max(Y_0,1)\) such that
\begin{equation}
 G_F(y)\leq C y^\theta\qquad(Y\leq y\leq Y^{1+\delta}).
 \label{eq:an-window}
\end{equation}
Then none of the \(L(s,\chi_i)\) has a zero with \(\Re s>\theta\).
The family, weights, exponent and constant are fixed across these
windows. Only their initial points need go to infinity.

\subsection{The transform and the sign of its residues}

Write \(g(u)=G_F(e^u)\) for \(u\geq0\). This is measurable. The individual
ramps are continuous across an integer because the new summand vanishes
there. Also \(|g(u)|\leq C_0e^{2u}\) for a constant depending on \(F\).
Indeed \(|\chi(n)|\leq1\) and \(\Lambda(n)\leq\log n\leq n\), so each
summand in \(B_\chi(y)\) has modulus at most \(y\), and there are at most
\(y\) summands. The centering terms have the same bound up to a fixed
constant. For \(\Re s>2\) consider the absolutely convergent transform
\[
 \Phi(s)=\int_0^\infty g(u)e^{-su}\,du .
\]
Changing variables to \(y=e^u\) gives the Mellin transform with measure
\(y^{-s-1}\,dy\). For \(n\geq1\),
\[
 \int_n^\infty (y/n-1)y^{-s-1}\,dy
       =\frac{n^{-s}}{s(s-1)} .
\]
The corresponding sum of absolute integrals converges because
\(\sum_{n\geq1}\Lambda(n)n^{-\Re s}\) converges for \(\Re s>1\).
Termwise integration followed by logarithmic differentiation of the
Euler product gives the following transforms \(\Phi_\chi\) of
\(G_\chi(e^u)\):
\begin{align*}
 \Phi_\chi(s)
 &=-\frac{L'(s,\chi)/L(s,\chi)}{s(s-1)}
          +\frac{L'(1,\chi)/L(1,\chi)}{s-1}
             &&(\chi\text{ nonprincipal}),\\
 \Phi_1(s)
 &=-\frac{\zeta'(s)/\zeta(s)}{s(s-1)}
          -\frac{1}{(s-1)^2}+\frac{1+\gamma}{s-1}.
\end{align*}
Here the transforms of \(y\) and \(y\log y\) are respectively
\((s-1)^{-1}\) and \((s-1)^{-2}\). The pole at one cancels in each
formula. For a nonprincipal character this follows by substitution
in the numerator after putting the terms over a common denominator.
For the principal character use
\(\zeta'(s)/\zeta(s)=-1/(s-1)+\gamma+O(s-1)\), which follows from the
Laurent expansion of \(\zeta\) at one. This is the divided-difference
removal of the singularity in the source.

On real \(u\) we have
\(\overline{G_\chi(e^u)}=G_{\chi^{-1}}(e^u)\). Indeed the Euler
coefficients are conjugate, and \(L(s,\chi^{-1})=
\overline{L(\overline{s},\chi)}\), first by the Dirichlet series and
then by continuation. Hence the meromorphic continuation of the transform is
\[
 \Phi(s)=\frac12\sum_i a_i\bigl(\Phi_{\chi_i}(s)
                                   +\Phi_{\chi_i^{-1}}(s)\bigr).
\]
Let \(Z\) be the set of zeros in \(\Re s>0\) of all the functions
\(L(s,\chi_i), L(s,\chi_i^{-1})\). Classical nonvanishing on
\(\Re s\geq1\) implies \(Z\subset\{0<\Re s<1\}\). A compact subset
of \(\Re s>0\) meets \(Z\) in finitely many points by isolatedness of
zeros and finiteness of the family. Away from \(Z\), the continuation is
holomorphic in that half-plane. For \(\rho\in Z\), write \(m_\chi(\rho)\)
for the zero multiplicity, taking it to be zero when the value is nonzero.
Locally \(L(s,\chi)=(s-\rho)^{m_\chi(\rho)}b(s)\) with \(b\)
holomorphic and \(b(\rho)\ne0\). Its logarithmic derivative is
\(m_\chi(\rho)/(s-\rho)+b'(s)/b(s)\). We therefore obtain
\begin{equation}
 \Phi(s)=\frac{c_\rho}{s-\rho}+\text{a holomorphic function},\qquad
 c_\rho=\frac{\sum_i a_i
             (m_{\chi_i}(\rho)+m_{\chi_i^{-1}}(\rho))}
                    {2\rho(1-\rho)} .
 \label{eq:an-res}
\end{equation}
The numerator is positive since at least one multiplicity is positive.
Thus \(c_\rho\ne0\). If \(\rho\) is real then \(0<\rho<1\), and
\(c_\rho\) is positive real. This last sign handles possible real zeros
in the same argument as the nonreal ones.

\subsection{What an upper bound on one window gives}

Take \(Y\) in \eqref{eq:an-window}, put \(U=\log Y\), and fix
\(\eta>0\). On \([U,(1+\delta)U]\) the function \(g\) has the nonnegative
majorant \(R(u)=|C|e^{\theta u}\). Its weighted integral is
\begin{equation}
 \int_U^{(1+\delta)U}R(u)e^{-\theta u}\,du
         =|C|\delta U\leq e^{\eta U}
 \label{eq:an-majorant}
\end{equation}
for all sufficiently large \(U\). We can take such \(U\) arbitrarily
large. The argument uses only \(g\leq R\), \(R\geq0\) and
\eqref{eq:an-majorant} on these intervals.

Put
\[
 A_k(z)=\frac{(-1)^k}{k!}\Phi^{(k)}(z)
       =\int_0^\infty g(u)\frac{u^k}{k!}e^{-zu}\,du
       \qquad(\Re z>2).
\]
For each fixed \(k\) and each closed neighborhood with
\(\Re z\geq2+b>2\), \(C_0u^ke^{-bu}\) is an integrable majorant;
the derivative formula follows by differentiation under the integral.
Since \(g\) is real, \(A_k(\sigma)\) is real for real \(\sigma>2\).
Taking absolute values on the vertical translate and adding the
untranslated integral yields, for every real \(t\),
\[
 |A_k(\sigma+it)|+A_k(\sigma)
  \leq 2\int_0^\infty \max(g(u),0)\frac{u^k}{k!}e^{-\sigma u}\,du .
\]
This follows pointwise from \(|g|+g=2\max(g,0)\). For \(u\geq0\),
the \(k\)-th term of the exponential series gives
\(u^ke^{-(\sigma-\theta)u}/k!\leq(\sigma-\theta)^{-k}\).
The contribution from the chosen interval is at most
\(2e^{\eta U}(\sigma-\theta)^{-k}\). Outside it we have the bound
\(2C_0u^ke^{-(\sigma-2)u}/k!\).

\paragraph{Concentration with the constants of the formal proof.}
For each \(\delta>0\) set
\[
 l=2+6/\delta,\quad M=l(1+\delta),\quad \sigma=M(2+l),\quad
 k=\left\lceil (1+\delta/2)\sigma U\right\rceil .
\]
Then \(\sigma>2\), and \(k\geq U\) when \(U\geq1\). Write
\(V=(1+\delta)U,\ a=\sigma-2-l,\ b=\sigma-2+l\); both \(a,b\) are positive.
The exact tail estimate used in the proof is
\begin{equation}
 \int_{(0,\infty)\setminus[U,V]}
 \frac{u^k}{k!}e^{-(\sigma-2)u}\,du
 \leq\left(\frac{e}{a}+\frac1b\right)\sigma^{-k}.
 \label{eq:an-tail}
\end{equation}
For \(u>V\) multiply the integrand by \(e^{l(u-V)}\geq1\), and for
\(u<U\) multiply it by \(e^{l(U-u)}\geq1\). Extending each resulting
nonnegative integral to \((0,\infty)\) gives
\(e^{-lV}a^{-k-1}+e^{lU}b^{-k-1}\), since repeated integration by parts
gives \(\int_0^\infty u^ke^{-zu}\,du=k!z^{-k-1}\) for \(z>0\).
For any \(z>0\), the inequality \(\log(\sigma/z)\leq(\sigma-z)/z\) yields
\[
 e^x z^{-k-1}\leq \frac{\sigma^{-k}}{z}
   \exp\!\left(x+k\frac{\sigma-z}{z}\right).
\]
For \(z=a\) the last ratio is \(1/(M-1)\). The identities
\((1+\delta/2)\sigma=M(M-1)\) and \(lV=MU\) show that
\(-lV+k/(M-1)\leq1\). For \(z=b\), the ratio is
\(-(l-2)/b\). Here \((1+\delta/2)(l-2)=l+1\) and \(lb\leq(l+1)\sigma\),
the latter following from \(l(l-2)\leq\sigma\). Thus \(k(l-2)/b\geq lU\),
so the second exponential is at most one. These two estimates prove
\eqref{eq:an-tail}.

Put \(E=\sigma-\theta\) and \(C_*=2+2C_0(e/a+1/b)\). We have proved
that for every \(\eta>0\) and every \(U_0\), some \(U\geq U_0\) and
integer \(k\geq U\) satisfy \emph{simultaneously for all} real \(t\)
\begin{equation}
 |A_k(\sigma+it)|+A_k(\sigma)\leq C_*e^{\eta U}E^{-k}.
 \label{eq:an-window-coeff}
\end{equation}
Indeed \(\sigma^{-k}\leq E^{-k}\) and we may choose \(U\geq1\).
The constants \(\sigma,C_*\) are independent of \(\eta,U_0,U,k,t\).

\subsection{Detecting a zero with high derivatives}

Suppose some \(\rho_0\in Z\) has \(\Re\rho_0>\theta\). The disk
\(\overline B(\sigma,E)\) is contained in the right half-plane since
\(\sigma-E=\theta>0\). If this disk contains a pole, let \(D\) be the
distance to the closest one; otherwise put \(D=E\). The open disk
\(B(\sigma,D)\) is pole-free.

First suppose some \(\rho_1\in Z\) satisfies \(\sigma-\Re\rho_1<D\).
At \(z_0=\sigma+i\Im\rho_1\) its distance is already \(<D\). Perturb
the imaginary part slightly, keeping this strict inequality, to obtain
\(z=\sigma+it\) with a unique nearest pole \(\rho\) and
\(d=|z-\rho|<D\). Here are the details. Choose \(\epsilon>0\) smaller
than both \(D-|z_0-\rho_1|\) and \(\sigma-D\). Only finitely many
poles lie in \(\overline B(z_0,D+\epsilon)\). For each pair from this
finite set with different imaginary parts, equality of squared distances
to \(\sigma+it\) is a linear equation in \(t\) with a nonzero coefficient.
For a pair of distinct poles with the same imaginary part, equality
would force equal real parts since both are to the left of \(\sigma\).
Choose \(t\) within \(\epsilon\) of \(\Im\rho_1\) outside the finite
exceptional set. All poles outside that disk have distance \(>D\)
from the chosen center. This proves the assertion.

Choose \(h\) strictly between \(d\) and the minimum of \(D\) and the
distance from \(z\) to any competing pole, using only poles in the
above compact disk. By \eqref{eq:an-res}, removing \(c_\rho/(s-\rho)\)
leaves a holomorphic function on a neighborhood of \(\overline B(z,h)\).
At the real center \(\sigma\), \(\Phi\) is holomorphic on a neighborhood
of the closed disk of radius \(h\), since \(h<D\). Cauchy's estimate on
these two circles gives \(M_1,M_2\), independent of \(k\), with
\[
 \left|A_k(z)-\frac{c_\rho}{(z-\rho)^{k+1}}\right|\leq M_1 h^{-k},
 \qquad |A_k(\sigma)|\leq M_2h^{-k}.
\]
Consequently the left side of \eqref{eq:an-window-coeff} is at
least \((|c_\rho|/d)d^{-k}-(M_1+M_2)h^{-k}\).

In the remaining case \(\sigma-\Re w\geq D\) for every pole \(w\).
The existence of \(\rho_0\) gives \(D<E\), so the defining disk
contains a closest pole. A pole at distance \(D\) has
\(D^2=(\sigma-\Re w)^2+(\Im w)^2\), forcing
\(w=\beta=\sigma-D\). Thus the closest pole is real and unique.
Take \(d=D\) and \(d<h<E\) small enough to exclude every other pole
from \(\overline B(\sigma,h)\). The same Cauchy estimate gives
\[
 A_k(\sigma)\geq(c_\beta/d)d^{-k}-M_1h^{-k}
\]
because \(c_\beta>0\) by \eqref{eq:an-res}. At \(t=0\) this also
bounds the left side of \eqref{eq:an-window-coeff} from below.

In either case there are \(c>0,\ 0<d<h\leq E\) and \(M'\geq0\)
such that \(c d^{-k}\leq M'h^{-k}+C_* e^{\eta U}E^{-k}\) on the
selected orders. Fix \(0<\eta<\log(h/d)\). Multiplication by \(d^k\)
and \(U\leq k\) bound the right side by
\(M'(d/h)^k+C_*(e^\eta d/h)^k\), which tends to zero as \(U_0\)
increases. This contradicts \(c>0\). All poles therefore satisfy
\(\Re\rho\leq\theta\). By construction each zero of a member of the
family with \(0<\Re s<1\) is one of these poles. The classical
nonvanishing theorem covers \(\Re s\geq1\), completing the window
criterion.

\subsection{The finite family attached to a primitive character}

Fix a primitive \(\psi\) of positive modulus \(N\). Let \(\Xi\) be the
finite cyclic group generated by \(\psi\) in the group of characters
modulo \(N\), and write \(o=|\Xi|\). Let \(S\) be the set of primitive
characters that induce the members of \(\Xi\), counted once.
It contains the modulus-one character, and it contains \(\psi\) because
\(\psi\) is already primitive. It is closed under inversion: induction
to a higher modulus commutes with inversion, and the conductors of
\(\chi\) and \(\chi^{-1}\) divide each other. Every principal member
of \(S\) therefore has modulus one, as also follows from the definition
of its conductor. Take the family indexed by \(S\), with weight \(1\)
on the principal character and \(1/(3o+1)\) on each other character.
For the cyclic level \(\mathcal L\) the formal construction chooses
no exceptional distinguished character. Its statistic is exactly
\[
 G_{\mathcal L}(y)=G_1(y)+\frac{1}{3o+1}
                  \sum_{\chi\in S\setminus\{1\}}\Re G_\chi(y)
                 =G_F(y).
\]
In this display \(G_1\) is real. Thus the preceding criterion applies
to every retained character with the arithmetic window held fixed.

Here is the precise arithmetic premise consumed by the analytic argument.
For this \(\mathcal L\) there exist \(e_D,e_L\in\mathbb R\) and
\(\varepsilon_1>0\) such that for every \(0<\varepsilon<\varepsilon_1\)
there are intervals \([Y,Y^{3/2}]\) with arbitrarily large \(Y\) on
which, simultaneously for all \(y\) in the interval,
\begin{equation}
 G_{\mathcal L}(y)\leq y^{1/2+(e_D+e_L+7)\varepsilon}.
 \label{eq:an-arith-input}
\end{equation}
In applying the formal comparison one first discards intervals lying
below the comparison threshold of the chosen family. The threshold
can depend on \(\varepsilon\), which is fixed before choosing \(Y\).

Take \(s\) with \(\Re s>1/2\), and put
\[
 t=\tfrac12\min(\Re s-1/2,1/2),\qquad
 \theta=1/2+t,\qquad
 \varepsilon=\min\!\left(\varepsilon_1/2,\,
                     \frac{t}{|e_D+e_L+7|+1}\right).
\]
Then \(1/2<\theta<1,\ \theta<\Re s\) and
\((e_D+e_L+7)\varepsilon\leq t\). On each of the selected windows
above one, \eqref{eq:an-arith-input} implies \eqref{eq:an-window} with
\(C=1,\delta=1/2\). The window criterion gives \(L(s,\psi)\ne0\),
subject to the pole convention already specified. This holds for
every primitive character because the construction can start with
any \(\psi\). The estimate \eqref{eq:an-arith-input} is the imported
arithmetic premise of this chapter; the analytic deduction does not
prove the arithmetic estimate.

\paragraph{Correspondence with the source.}
The ramp, positivity and admissibility conditions are the definitions
in \texttt{Deduction.\allowbreak{}Char\allowbreak{}Family}. The transform and residue calculation
translate \texttt{Deduction.\allowbreak{}pole\allowbreak{}Data}. The moment estimate and geometry
of the nearest pole translate \texttt{landau\_\allowbreak{}tails},
\texttt{landau\_\allowbreak{}window\_\allowbreak{}key}, \texttt{landau\_\allowbreak{}dichotomy},
\texttt{landau\_\allowbreak{}nearest} and \texttt{landau\_\allowbreak{}window}. The transition
between window hypotheses and zero-freeness is \texttt{lemma\allowbreak{}L} and
\texttt{zero\allowbreak{}Free\_\allowbreak{}of\_\allowbreak{}rpow\_\allowbreak{}window}. The cyclic family and the choice of
the exponent translate \texttt{Upper\allowbreak{}Theory.\allowbreak{}Comparison\ensuremath{{}_2}.\allowbreak{}level\_\allowbreak{}family},
\texttt{form\allowbreak{}B}, \texttt{Formalism.\allowbreak{}HC.\allowbreak{}level\_\allowbreak{}of\_\allowbreak{}tied\allowbreak{}Win} and
\texttt{zero\allowbreak{}Free\allowbreak{}On\_\allowbreak{}of\_\allowbreak{}tied\allowbreak{}Win\_\allowbreak{}cyc}. Arithmetic supplies the paired
inputs to the last theorem through \texttt{level\_\allowbreak{}gm\_\allowbreak{}le\_\allowbreak{}on\_\allowbreak{}family\_\allowbreak{}scales}
and \texttt{Upper\allowbreak{}Theory.\allowbreak{}Top.\allowbreak{}zero\allowbreak{}Free\allowbreak{}All\_\allowbreak{}of\_\allowbreak{}top\allowbreak{}G}.
